\BourbakiExerciseGroup

\begin{enumerate}
\def\labelenumi{\arabic{enumi})}
\item
  设 A 是一个交换环。证明：如果每个自由 A-模的每个子模都是自由的，那么 A 的每个理想都是主理想，并且 A 是整环（从而是一个主理想整环）。
\item
  设 M 是主理想整环 A 上的一个模；假设 M 是循环子模的直和。证明 M 的每个子模 N 都是循环子模的直和（若 T 是 N 的挠子模，首先利用定理 1 证明 N/T 是自由模，并由此推出 T 在 N 中有补；然后利用第七章第 57 页习题 12 c) 证明 T 是循环子模的直和）。
\item
  设 \((p_n)\) 是一个严格递增的素数序列，使得 \(p_n\) 不整除下列数中的任何一个
\end{enumerate}

\[
a + b \left(1 + p _ {1} + p _ {1} p _ {2} + \dots + p _ {1} p _ {2} \dots p _ {n - 1}\right) \quad \text { for } \quad | a | \leqslant n \quad \text { and } \quad | b | \leqslant n.
\]

在有理平面 \(\mathbf{Q}^2\) 中，考虑 \(\mathbf{Z}\) -模 \(\mathbf{M}\) ，它由如下定义的序列 \((x_n)\) 生成：

\[
x _ {0} = (1, 0), \quad x _ {1} = (0, 1), \quad x _ {n + 1} = p _ {n} ^ {- 1} \left(x _ {0} + x _ {n}\right) \quad \text { for } \quad n \geqslant 1.
\]

证明 M 的秩为 2，并且 M 的每个秩 1 子模都是循环的（注意 \(x_{0}\) 和 \(x_{n}\) 构成 M 的由诸 \(x_{i}\) （满足 \(i \leqslant n\) ）生成的子模的一组基）。

\begin{enumerate}
\def\labelenumi{\arabic{enumi})}
\setcounter{enumi}{3}
\tightlist
\item
  设 \(E\) 是主理想整环 \(A\) 上的无挠模。
\end{enumerate}

\begin{enumerate}
\def\labelenumi{\alph{enumi})}
\item
  证明：若F是E的纯子模（VII，第55页，习题7），则E/F是无挠的；反之亦然。
\item
  设 F 为 E 的由一个极大拟自由族（VII，第 55 页，习题 4）生成的子模；则 F 是 E 的纯的自由子模。证明：对所有 \(\dot{x} \in \mathrm{E} / \mathrm{F}\) ，存在一个属于 A 的不可逆元 \(\alpha\) 以及一个元素 \(\dot{y} \in \mathrm{E} / \mathrm{F}\) ，使得 \(\dot{x} = \alpha \dot{y}\) ，但 E/F 不一定是可除模（参见习题 3）。
\end{enumerate}

\begin{enumerate}
\def\labelenumi{\arabic{enumi})}
\setcounter{enumi}{4}
\tightlist
\item
  设 A 是主理想整环，且具有唯一的极大理想 \(\mathbf{A}\pi\) （参见 VII，第 48 页，习题 4）。设 E 是无挠 A-模，M 是 E 的自由子模，使得 E/M 是无挠、可除且秩为 1 的。设 \((e_{\alpha})\) 是 M 的一组基，设 \(\dot{u}\) 是 E/M 的一个非零元素，且 \(u\) 是陪集 \(\dot{u}\) 的一个元素；对每个整数 \(n > 0\) ，设 \(u_{n}\) 是 E 的一个元素，使得 \(u \equiv \pi^{n} u_{n} (\bmod M)\) 。令 \(u = \pi^{n} u_{n} + \sum_{\alpha} \lambda_{n\alpha} e_{\alpha}\) 。
\end{enumerate}

\begin{enumerate}
\def\labelenumi{\alph{enumi})}
\item
  证明 M 由 M 和诸 \(u_{n}\) 生成，并且 \(\lambda_{n\alpha} - \lambda_{n + 1,\alpha} \in \mathbf{A}\pi^n\) 对所有 \(\alpha\) 并且对每个整数 \(n > 0\) 成立（参见习题 4，a)）。
\item
  进一步假设 A 关于拓扑 T 是完备的，T 由取理想 \(A\pi^{n}\) 作为 A 中 0 的开邻域基而定义（一般拓扑，III，第 5 页）；令 \(\lambda_{\alpha} = \lim_{n \to \infty} \lambda_{n\alpha}\) 。证明：如果 \(\mu\) 和 \(\mu_{\alpha}\) 使得 \(\mu u - \sum_{\alpha} \mu_{\alpha} e_{\alpha}\) 属于所有
\end{enumerate}

子模 \(\pi^n\mathrm{E}\) ，则必然有 \(\mu_{\alpha} = \lambda_{\alpha}\mu\) 对所有 \(\alpha\) 。

\begin{enumerate}
\def\labelenumi{\arabic{enumi})}
\setcounter{enumi}{5}
\tightlist
\item
  设 \(\mathbf{A}\) 是一个主理想整环，具有唯一的极大理想 \(\mathbf{A}\pi\) ，并且关于习题 5 中定义的拓扑 \(\mathcal{T}\) 是完备的。
\end{enumerate}

\begin{enumerate}
\def\labelenumi{\alph{enumi})}
\item
  证明：任意秩有限的、无挠的 A-模 E，若不包含非零可除子模，则 E 是自由的（考虑由一个极大拟自由族生成的 E 的子模 M；若 \(M \neq E\) 则借助习题 4，b) 和 5，b) 证明 E 包含非零可除子模）。将此结果推广到 E 具有可数秩的情形（用归纳法进行）。
\item
  在自由模 \(E = A^{(N)}\) （A 的可数无穷多个副本的直和）中，设 \((a_{n})_{n \geq 0}\) 为典范基，并令 \(e_{n} = a_{n-1} - \pi a_{n}\) 对每个整数 \(n \geq 1\) 。证明由诸 \(e_{n}\) 生成的 E 的子模 M 是纯的，并且诸 \(e_{n}\) 构成一个极大的拟自由族，但 E/M 是可除模，因此 M 在 E 中没有补。
\end{enumerate}

\begin{enumerate}
\def\labelenumi{\arabic{enumi})}
\setcounter{enumi}{6}
\tightlist
\item
  设 \(p\) 是一个素数，并设 \(\mathbf{A}\) 表示形如 \(r / s\) 的有理数组成的环，其中 \(s\) 与 \(p\) 互素（VII，第 48 页，习题 4）。
\end{enumerate}

设 \(u = (1,0)\) 和 \(v = (0,1)\) 是向量空间 \(E = Q^2\) （ \(Q\) 上的）的典范基的元素。递归地定义一个元素序列 \(u_n\) （ \(E\) 中的）： \(u_0 = u\) ，且 \(u_{n}=pu_{n+1}+v\) 对 \(n\geqslant0\) 。设 F 为由 v 和诸 \(u_{n}\) 生成的 A-子模；证明 F 不是自由模，但不含非零可除子模（参见习题 6，a)）。

\begin{enumerate}
\def\labelenumi{\arabic{enumi})}
\setcounter{enumi}{7}
\tightlist
\item
  设 A 是一个非域的主理想整环。
\end{enumerate}

\begin{enumerate}
\def\labelenumi{\alph{enumi})}
\item
  证明乘积 A-模 \(A^{N}\) 没有可数生成系。（若 A 可数，考虑 \(\operatorname{Card}(A^{N})\) ；否则注意到，若对所有 \(x \in A\) 我们令 \(v_{x}\) 表示元素 \((x^{n})_{n \geq 0}\) （属于 \(A^{N}\) ），则诸 \(v_{x}\) 构成自由系）。
\item
  设 \(\pi\) 为 A 的一个不可约元素。令 S 表示 \(\mathbf{A}^{\mathbf{N}}\) 的子模，它由这样的序列 \(z = (z_{n})_{n\in \mathbb{N}}\) 组成：存在一个整数序列 \((k(n))_{n\in \mathbb{N}}\) ，它趋于 \(+\infty\) （依赖于 \(z\) ），使得 \(z_{n}\in \mathbf{A}\pi^{k(n)}\) 对所有 \(n\) 成立。证明 \((\mathbf{A} / \pi)\) -向量空间 \(\mathrm{S} / \pi \mathrm{S}\) 有一个可数基。
\item
  由 a) 和 b) 推出， \(A^{N}\) 不是自由 A-模。（否则，S 将是一个具有可数基的自由 A-模；但 \((x_{n})\mapsto(x_{n}\pi^{n})\) 是一个从 \(A^{N}\) 到 S 中的单射 A-线性映射。）
\end{enumerate}

\begin{enumerate}
\def\labelenumi{\arabic{enumi})}
\setcounter{enumi}{8}
\tightlist
\item
  设 \(\mathbf{A}\) 为主理想整环，且设 \(\mathbf{M}\) 为 \(\mathbf{A}\) -模 \(\mathbf{A}^{\mathbf{N}}\) 的对偶模。设 \(f: \mathbf{M} \to \mathbf{A}^{\mathbf{N}}\) 为典范单射 \(\mathbf{A}^{(\mathbf{N})} \to \mathbf{A}^{\mathbf{N}}\) 的转置映射（其中 \(\mathbf{A}^{\mathbf{N}}\) 与 \(\mathbf{A}^{(\mathbf{N})}\) 的对偶等同）。
\end{enumerate}

\begin{enumerate}
\def\labelenumi{\alph{enumi})}
\item
  设 \(\pi\) 为 A 的一个不可约元素，且令 \(\varphi: \mathbf{A}^{\mathbf{N}} / \mathbf{A}^{(\mathbf{N})} \to \mathbf{A}\) 是一个线性形式；证明如果 \(u \in \mathbf{A}^{\mathbf{N}} / \mathbf{A}^{(\mathbf{N})}\) 是某个元素 \((u_n) \in \mathbf{A}^{\mathbf{N}}\) （满足 \(u_n \in (\pi^n)\) 对所有 \(n\) ）的类，则 \(\varphi(u) = 0\) 。如果 A 包含两个非相伴的不可约元素，则推出 \(f\) 是单射的。（利用上述结论和贝祖恒等式，证明 \(\mathbf{A}^{\mathbf{N}} / \mathbf{A}^{(\mathbf{N})}\) 的对偶为零。）
\item
  给 A 赋予 VII 第 60 页习题 5 中所定义的拓扑。证明：若 \(f(\mathbf{M})\) 不包含于 \(\mathbf{A}^{(\mathbf{N})}\) ，则 A 是完备的。（将问题归结为证明：若 \(\varphi : \mathbf{A}^{\mathbf{N}} \to \mathbf{A}\) 是一个线性形式，使得 \(\varphi(e_n) \neq 0\) 对每个元素 \(e_n\) （属于 \(\mathbf{A}^{(\mathbf{N})}\) 的典范基）成立，则 \(\operatorname{Im}(\varphi)\) 包含 A 中任意一个趋于 0 足够快的元素级数的和。）
\item
  如果 A 包含两个非相伴的不可约元素，则从 \(\mathbf{A}^{(N)}\) （相应地 \(A^{N}\) ）到其双重对偶的典范映射是双射（利用 a) 和 b)）。
\end{enumerate}

\begin{enumerate}
\def\labelenumi{\arabic{enumi})}
\setcounter{enumi}{9}
\item
  设 M 为 p 进整数环 \(Z_{p}\) 的一个非零 Z-子模，且是 Z-模 \(Z_{p}\) （VII，第 55 页，习题 7）的纯 Z-子模。证明 \(M + pZ_{p} = Z_{p}\) ；由此推出 M 是不可分解的（注意 \(Z_{p}/pZ_{p}\) 同构于 Z/pZ）。给出 \(Z_{p}\) 的一个可分解 Z-子模的例子（观察到 \(Z_{p}\) 具有连续统的基数）。
\item
  设 A 和 B 是两个无限的有理素数集合，它们没有公共元素，且都不包含数 5。设 \((e_i)_{1 \leqslant i \leqslant 4}\) 是 Q-向量空间 \(\mathbf{Q}^4\) 的典范基。
\end{enumerate}

\begin{enumerate}
\def\labelenumi{\alph{enumi})}
\item
  设 M 为 \(Q^{4}\) 的由元素 \(e_{1}/m\) （给定 \(m \in A\) ）生成的 Z-子模；设 N 为 \(Q^{4}\) 的由元素 \(e_{2}/m\) （其中 \(m \in A\) ）、 \(e_{3}/n\) （其中 \(n \in B\) ）以及 \((e_{2} + e_{3})/5\) 生成的 Z-子模，则 \(M \cap N = 0\) 。令 \(e_{1}' = 8e_{1} + 3e_{2}\) 和 \(e_{2}' = 5e_{1} + 2e_{2}\) 。设 \(M'\) 为 \(Q^{4}\) 的由元素 \(e_{1}'/m\) （其中 \(m \in A\) ）生成的 Z-子模；设 \(N'\) 为 \(Q^{4}\) 的由元素 \(e_{2}'/m\) （其中 \(m \in A\) ）、 \(e_{4}/n\) （其中 \(n \in B\) ）以及 \((3e_{2}' + e_{3})/5\) 生成的 Z-子模。证明 \(M' \cap N' = 0\) ， \(M \oplus N = M' \oplus N'\) ；M 与 \(M'\) 同构且不可分解；并且 N 与 N' 不可分解但不同构（比较 VII，第 69 页，习题 23）。
\item
  设 M 为 \(Q^{4}\) 的由元素 \(e_{1}/m\) （其中 \(m \in A\) ）、 \(e_{2}/n\) （其中 \(n \in B\) ）以及 \((e_{1} + e_{2})/5\) 生成的 Z-子模；设 N 为 \(Q^{4}\) 的由元素 \(e_{3}/m\) （其中 \(m \in A\) ）、 \(e_{4}/n\) （其中 \(n \in B\) ）以及 \((e_{3} + e_{4})/5\) 生成的 Z-子模。令 \(e_{1}' = 2e_{1} + e_{3}\) , \(e_{2}' = e_{2} + 3e_{4}\) , \(e_{3}' = 17e_{1} + 9e_{3}\) ，以及 \(e_{4}' = e_{2} + 2e_{4}\) 。设 \(M'\) 为 \(Q^{4}\) 的由元素 \(e_{1}'/m\) （给定 \(m \in A\) ）、 \(e_{2}'/n\) （给定 \(n \in B\) ）以及 \((e_{1}' + 2e_{2}')/5\) 生成的 Z-子模；设 \(N'\) 为 \(Q^{4}\) 的由元素 \(e_{3}'/m\) （给定 \(m \in A\) ）、 \(e_{4}'/n\) （给定 \(n \in B\) ）以及 \((e_{3}' + 2e_{4}')/5\) 生成的 Z-子模。证明 \(M \oplus N = M' \oplus N'\) ，M 与 N 同构，且 \(M'\) 同构于 \(N'\) ，但 \(M'\) 与 M 不同构。
\item
  设 M 为 \(Q^{4}\) 的由 \(e_{1}/5^{n}\) ( \(n \geqslant 0\) ) 生成的 Z-子模；设 N 为 \(Q^{4}\) 的由 \(e_{2}/5^{n}\) , \(e_{3}/2^{n}\) , \(e_{4}/3^{n}\) ( \(n \geqslant 0\) ), \((e_{2} + e_{3})/3\) 和 \((e_{2} + e_{4})/2\) 生成的子模。令 \(e_{1}' = 3e_{1} - e_{2}\) 和 \(e_{2}' = 2e_{1} - e_{2}\) 。设 \(M'\) 为 \(Q^{4}\) 的由元素 \(e_{1}'/5^{n}\) , \(e_{3}/2^{n}\) ( \(n \geqslant 0\) ）和 \((e_{1}' - e_{3})/8\) 生成的 Z-子模；设 \(N'\) 为 \(Q^{4}\) 的由元素 \(e_{2}'/5^{n}\) , \(e_{4}/3^{n}\) ( \(n \geqslant 0\) ）和 \((e_{2}' - e_{4})/2\) 生成的 Z-子模。证明 \(M \oplus N = M' \oplus N'\) ， \(M'\) 和 \(N'\) 是秩为 2 的不可分解模，且 M 的秩为 1。
\end{enumerate}
% semantic-fragments: legacy-input-seam
\relax{}
